\chapter{Autenticación con Google Sign-In}
\label{ch:google_signin}

\section{Objetivo}

Implementar el inicio de sesión con Google, restringido exclusivamente a cuentas institucionales \texttt{@cua.uam.mx}, permitiendo a los usuarios acceder a la plataforma con un solo clic, sin necesidad de recordar contraseñas.

\section{Contexto}

Dado que la UAM-C utiliza Google Workspace for Education, todos los correos \texttt{@cua.uam.mx} son cuentas de Google. Integrar Google Sign-In no solo mejora la experiencia de usuario, sino que refuerza la seguridad, ya que delega la autenticación a Google y valida el dominio institucional.

\section{Desarrollo}

\subsection{Configuración en Google Cloud Console}

Se generó un \textbf{ID de cliente de OAuth} de tipo \textit{Android}. Se registró el nombre del paquete \texttt{mx.uam.notiuam\_mobile} y la huella digital SHA-1 del emulador.

Para obtener el SHA-1 del entorno de desarrollo, se ejecutó:

\begin{lstlisting}[language=bash]
cd android && ./gradlew signingReport
\end{lstlisting}

\subsection{Backend}

Se creó el servicio \texttt{GoogleAuthService} que valida el ID Token enviado por el cliente. Utiliza la librería de Google API Client y verifica que la audiencia coincida con el Client ID configurado.

\begin{lstlisting}[language=Java]
@Service
public class GoogleAuthService {
    @Value("${GMAIL_CLIENT_ID}")
    private String clientId;

    public GoogleIdToken.Payload verify(String idTokenString) {
        GoogleIdTokenVerifier verifier = new GoogleIdTokenVerifier.Builder(
                new NetHttpTransport(), GsonFactory.getDefaultInstance())
                .setAudience(Collections.singletonList(clientId))
                .build();
        GoogleIdToken idToken = verifier.verify(idTokenString);
        if (idToken == null) throw new IllegalArgumentException("Token invalido");
        return idToken.getPayload();
    }
}
\end{lstlisting}

En \texttt{AuthService} se añadió el método \texttt{loginWithGoogle} que recibe el token, valida el dominio y crea el usuario si no existe (marcándolo como verificado directamente):

\begin{lstlisting}[language=Java]
public AuthResponse loginWithGoogle(String idTokenString) {
    var payload = googleAuthService.verify(idTokenString);
    String email = payload.getEmail().toLowerCase().trim();

    if (!email.endsWith("@cua.uam.mx")) {
        throw new IllegalArgumentException("Solo se permiten correos @cua.uam.mx");
    }

    User user = userRepository.findByEmail(email).orElseGet(() -> {
        User nuevo = User.builder()
                .email(email)
                .displayName(payload.get("name") != null ? payload.get("name").toString() : email)
                .passwordHash("") // Sin contrasena, solo Google
                .role(Role.STUDENT)
                .emailVerified(true) // Google ya verifica
                .enabled(true)
                .build();
        return userRepository.save(nuevo);
    });

    return buildAuthResponse(user);
}
\end{lstlisting}

\subsection{Cliente Flutter}

Se añadió la dependencia \texttt{google\_sign\_in}. El servicio \texttt{GoogleAuthService} utiliza el parámetro \texttt{hostedDomain} para restringir el selector de cuentas a solo \texttt{@cua.uam.mx}.

\begin{lstlisting}[language=Java]
class GoogleAuthService {
  static final _googleSignIn = GoogleSignIn(
    scopes: ['email', 'profile'],
    hostedDomain: 'cua.uam.mx',
    serverClientId: '584002587099-...apps.googleusercontent.com',
  );

  static Future<String?> signInAndGetIdToken() async {
    final account = await _googleSignIn.signIn();
    if (account == null) return null;
    final auth = await account.authentication;
    return auth.idToken;
  }
}
\end{lstlisting}

Se añadió un botón de \textit{Iniciar sesión con Google} en la pantalla de login, con un separador visual que divide el formulario tradicional del inicio con Google.

\subsection{Problemas encontrados y solución}

\subsubsection{Error \texttt{ApiException: 10} en Android}

Al intentar el login, el emulador lanzaba \texttt{ApiException: 10 (DEVELOPER\_ERROR)}. \textbf{Solución:} 
\begin{enumerate}[leftmargin=*]
  \item Se agregó el SHA-1 del emulador en Firebase Console.
  \item Se descargó el nuevo \texttt{google-services.json}.
  \item Se detectó que faltaba el \texttt{client\_type: 1} (Android) en el archivo, por lo que se creó manualmente el Client ID de Android en Google Cloud Console.
  \item Se forzó el \texttt{serverClientId} en el código de Flutter con el Client ID de tipo Web.
\end{enumerate}

\section{Verificación}

\begin{enumerate}[leftmargin=*]
  \item Abrir la app en el emulador.
  \item Tocar \textbf{Iniciar sesión con Google}.
  \item El selector de cuentas solo muestra correos \texttt{@cua.uam.mx}.
  \item Seleccionar la cuenta y entrar directamente al \texttt{HomeShell}.
\end{enumerate}

\begin{exito}[Resultado obtenido]
El flujo de Google Sign-In funciona correctamente y el usuario accede a la app sin ingresar contraseña.
\end{exito}

\section{Resultado}

\begin{exito}[Estado al cierre]
\begin{itemize}[leftmargin=*]
  \item Autenticación con Google restringida a cuentas institucionales.
  \item Creación automática de usuarios si es su primer inicio de sesión.
  \item Eliminación de la fricción del registro con contraseña para usuarios de Google Workspace.
\end{itemize}
\end{exito}
